\section{A Checked and Modular Proof Workflow}
\label{sec:proof-safety}

The encoding in \cref{sec:partial-operators} requires explicit WD evidence, and Lean's kernel checks each theorem application.
We examine two complementary aspects of proof construction through examples in Atelier~B\footnote{Community Edition 24.04.2.}: the definedness of terms introduced during interaction and the validity of the inference rules applied.

\subsection{Well-definedness throughout interaction}
\label{sec:lean-ip-comparison}

The WD obligations generated from a B specification concern the expressions occurring in that specification.
Interactive proof can introduce further partial expressions, for instance when choosing an existential witness.
Their definedness must be justified at that point.

Atelier~B's reference manual states that the tactic \texttt{se}, which supplies an existential witness, does not protect against ill-typing or ill-definedness~\cite[\S4.44]{AtelierBInteractiveProver}.
Consider the machine defined in \texttt{specs/DerivFalse.mch} containing the contradictory assertion \(\exists x \in \mathbb{Z}, x \notin \mathbb{Z}\).
After an initial deduction step (\texttt{dd}), the command \inlineab!se(min({}))! selects the minimum of the empty set as witness.
The displayed goal becomes
\[
  \bmin{\varnothing} \in \mathbb{Z}
  \;\land\;
  \bmin{\varnothing} \notin \mathbb{Z}.
\]
A call to \texttt{mp}, the mini-prover, closes it.
In this example, the prover accepts a contradictory assertion after an undefined witness has been supplied interactively.

In \barrel, forming \(\min S\) requires evidence that $S$ has a least element (\cref{lst:min-fn-app}).
The witness \(\min\varnothing\) thus requires a proof of \(\mathtt{min.WD}\ \varnothing\), that is,
\[
  \exists x \in \varnothing,\ \forall y \in \varnothing,\ x \le y,
\]
which is unprovable.
Admitting this WD condition triggers a Lean warning, and the \texttt{assert\_no\_sorry} check in \texttt{examples/DerivFalse.lean} fails on the resulting theorem.
For operators encoded with WD parameters, definedness is therefore checked throughout proof construction, including steps the source PO generator could not anticipate.

\subsection{Validity of inference rules}
\label{sec:checked-rules}

Inference rule validity is a separate requirement from definedness.
The rule \texttt{GenPredicateX.8} in Atelier~B's supplied database has the following schema:
\[
  \frac{a \notin \mathcal{P}(X) \cup \{y\}}
       {\exists z\in a,\ z\notin X \land z\ne y}.
\]
The premise places $y$ at the type of $a$, whereas the conclusion compares it with an element of $a$.
The schema therefore has no consistent B typing for arbitrary \(y\).
Instantiating $y$ with the polymorphic empty set nevertheless gives individually well-typed formulas and an invalid inference.

Indeed, let $a := \{\varnothing\}$ and $X := \mathcal{P}_1(\mathbb{Z})$, the set of nonempty subsets of $\mathbb{Z}$.
Since $\varnothing \in a$ but $\varnothing \notin X$, we have $a \notin \mathcal{P}(X)$. Since $a \ne \varnothing$,
the premise $a\notin\mathcal{P}(X)\cup\{\varnothing\}$ holds.
The conclusion requires an element of $\{\varnothing\}$ distinct from $\varnothing$, which is impossible.
\looseness=-1


The following script derives this contradiction and closes the current goal using only total operators:
\begin{lstlisting}[language=b, basicstyle=\small\ttfamily, numbers=none, mathescape=false, escapeinside={}]
ah(not({{}}: POW(POW1(INTEGER)) \/ {{}})) & pr & dd &
  ar(GenPredicateX.8,Fwd) & pr
\end{lstlisting}

In Lean, we refute this rule instance and prove that replacing \(\mathcal{P}(X) \cup\{y\}\) by $\mathcal{P}(X\cup\{y\})$ in the premise yields a valid rule.
Both proofs are available in \texttt{examples/Falsehood.lean} in the artifacts.

These counterexamples motivate two requirements for a proof workflow: terms introduced during interaction must carry the definedness evidence required by their operators, and reusable inference rules must themselves have checked proofs.
The case study in \cref{sec:case-study} examines how these requirements interact with lemma reuse, refinement, and automation in a complete B development.
